% ============ §5 指示代词与 нравиться ============
\hdr{§5}{指示代词、最高级与 нра́виться：把「这个」和「喜欢」说对}

\callout{\textbf{这一节的坐标}：У8 p2--p3（э́тот／э́та／э́то／э́ти）、p5--p8（са́мый 与四个对比对话）、p12--p19（нра́виться ＋ 第三格）、p20--p23（знать／ду́мать, что ＋ потому́ что）、p27--p28（ИК-5）。\\
\textbf{这一节要解决的事}：看到一个名词，能选出 э́тот 的哪一个形式；说「我喜欢」时，知道谁做主语、谁变第三格。}

\subhdr{5.1　э́тот／э́та／э́то／э́ти：跟 како́й 一一对应}

У8 p3 把两组东西并排放在一起：左边是 како́й 的四个疑问形式，右边是「这个」的四个形式。它们是一一对应的——\textbf{问句用哪个词尾，答句就用哪个}：

\begin{minipage}{\linewidth}
\tabtop
\begin{tabular}{@{}>{\raggedright\arraybackslash}p{0.22\linewidth}>{\raggedright\arraybackslash}p{0.22\linewidth}>{\raggedright\arraybackslash}p{0.42\linewidth}@{}}
\thead{疑问词} & \thead{指示代词} & \thead{课件例（У8 p5--p8）} \\
\end{tabular}
\tabmid
\begin{tabular}{@{}>{\raggedright\arraybackslash}p{0.22\linewidth}>{\raggedright\arraybackslash}p{0.22\linewidth}>{\raggedright\arraybackslash}p{0.42\linewidth}@{}}
Како́й …? & э́тот & Вот э́тот телефо́н. \\
Кака́я …? & э́та & Вот э́та ма́ленькая матрёшка. \\
Како́е …? & э́то & Вот э́то бе́лое пальто́. \\
Каки́е …? & э́ти & Вот э́ти часы́, спра́ва. \\
\end{tabular}
\tabbot
\end{minipage}

\textbf{这张表在说什么}：四个形式的分工与 \S4.2 那五族完全一样——跟名词的性数走。最容易错的是第三行：\textbf{中性用 э́то，不是 *э́тот}。课件 У8 p2 的中性例句写的就是 \textbf{Э́то пальто́ са́мое тёплое.}

\textbf{逐行读法}：第一行 телефо́н 阳性 → э́тот；第二行 матрёшка 阴性 → э́та；第三行 пальто́ 中性 → э́то；第四行 часы́ 只有复数 → э́ти。

\warnbox{\textbf{两个「э́то」不是一回事}。问句 \textbf{Что э́то?} 里的 э́то 是「这（是）」，\textbf{不修饰名词、不跟性数变}；而 э́тот／э́та／э́то／э́ти 是\textbf{指示代词}，必须修饰名词、必须跟性数。所以「这是书吗」用 \textbf{Э́то кни́га?}，而「这本书」用 \textbf{э́та кни́га}——前者的 э́то 不跟着 кни́га 变成 э́та。}

\subhdr{5.2　最高级 са́мый：形容词也要跟着名词变}

У8 p2 与 p5--p8 的四个对话是同一个模板：问 \textbf{Како́й … са́мый …?}，答 \textbf{Вот э́тот … са́мый …}。最高级不是加一个词就完了，\textbf{са́мый 和后面的形容词都要跟名词一致}：

\begin{minipage}{\linewidth}
\tabtop
\begin{tabular}{@{}>{\raggedright\arraybackslash}p{0.10\linewidth}>{\raggedright\arraybackslash}p{0.16\linewidth}>{\raggedright\arraybackslash}p{0.56\linewidth}@{}}
\thead{性数} & \thead{са́мый 的形式} & \thead{课件例句} \\
\end{tabular}
\tabmid
\begin{tabular}{@{}>{\raggedright\arraybackslash}p{0.10\linewidth}>{\raggedright\arraybackslash}p{0.16\linewidth}>{\raggedright\arraybackslash}p{0.56\linewidth}@{}}
阳 & са́мый & Како́й телефо́н са́мый дорого́й? — Э́тот телефо́н са́мый дорого́й. \\
阴 & са́мая & Кака́я матрёшка са́мая дешёвая? — Вот э́та матрёшка са́мая дешёвая. \\
中 & са́мое & Како́е пальто́ са́мое тёплое? — Вот э́то пальто́ о́чень тёплое. \\
复 & са́мые & Каки́е часы́ са́мые хоро́шие? — Вот э́ти часы́. \\
\end{tabular}
\tabbot
\end{minipage}

\textbf{这张表在说什么}：一句话——\textbf{са́мый 是「最」，它自己是一个形容词，所以要和后面的形容词一起变}。两处都变对，句子才对：\textbf{са́мая краси́вая}（阴性两个 -ая）、\textbf{са́мое тёплое}（中性两个 -ое）。

\textbf{逐行读法}：第三行的答案里课件没有重复 са́мое тёплое，只说 \textbf{о́чень тёплое}——这是口语省略，完整形式仍是 \textbf{са́мое тёплое}。写作业时建议写完整。

\subhdr{5.3　нра́виться：主语是「被喜欢的东西」，喜欢的人是第三格}

У8 p12--p18 的六个对话都在做同一件事。先看课件给的第三格表（У8 p14）：

\begin{minipage}{\linewidth}
\tabtop
\begin{tabular}{@{}>{\raggedright\arraybackslash}p{0.10\linewidth}>{\raggedright\arraybackslash}p{0.12\linewidth}>{\raggedright\arraybackslash}p{0.18\linewidth}>{\raggedright\arraybackslash}p{0.10\linewidth}>{\raggedright\arraybackslash}p{0.12\linewidth}>{\raggedright\arraybackslash}p{0.18\linewidth}@{}}
\thead{谁} & \thead{第三格} & \thead{例} & \thead{谁} & \thead{第三格} & \thead{例} \\
\end{tabular}
\tabmid
\begin{tabular}{@{}>{\raggedright\arraybackslash}p{0.10\linewidth}>{\raggedright\arraybackslash}p{0.12\linewidth}>{\raggedright\arraybackslash}p{0.18\linewidth}>{\raggedright\arraybackslash}p{0.10\linewidth}>{\raggedright\arraybackslash}p{0.12\linewidth}>{\raggedright\arraybackslash}p{0.18\linewidth}@{}}
я & мне & Мне нра́вится фи́зика. & мы & нам & Нам нра́вится ру́сский язы́к. \\
ты & тебе́ & Тебе́ нра́вится э́тот парк. & вы & вам & Вам нра́вится …? \\
он & ему́ & Ему́ нра́вится го́род. & они́ & им & Им нра́вится зима́. \\
она́ & ей & Ей нра́вится спорт. & — & — & — \\
\end{tabular}
\tabbot
\end{minipage}

\textbf{这张表在说什么}：它只给了「谁变成什么」，没给「为什么变」。原因是 нра́виться 的句法结构跟汉语相反：\textbf{句子的主语不是「喜欢的人」，而是「被喜欢的东西」}——所以那个东西决定动词用 \textbf{нра́вится}（单数）还是 \textbf{нра́вятся}（复数），而喜欢的人只能用第三格。

\textbf{逐行读法}：\textbf{Мне нра́вится фи́зика.} 里主语是 фи́зика（单数）→ нра́вится；\textbf{Мне нра́вятся кита́йские фи́льмы.} 里主语是 фи́льмы（复数）→ нра́вятся。人（мне）始终不变格——它已经是第三格了。

\aside{这条结构在 \S3 已经见过一半：第三格 мне／тебе́／ему́／ей／нам／вам／им 就是「谁」这一列的变格。这里的新东西是\textbf{动词跟谁一致}——跟被喜欢的东西，不跟人。}

\warnbox{\textbf{边界一}：нра́виться 的「被喜欢的东西」是主语，所以要用第一格（фи́зика、фи́льмы），不能用第四格。\\
\textbf{边界二}：\textbf{Мне нра́вится} 里的 мне 不能换成 я——\textbf{*Я нра́вится} 是初学者最常见的错句。}

\subhdr{5.4　знать／ду́мать, что … 与 потому́ что}

У8 p20--p23 给了两个动词的变位和两个复合句结构：

\begin{itemize}
\item \textbf{знать}（知道）：я зна́ю, ты зна́ешь, он зна́ет, мы зна́ем, вы зна́ете, они́ зна́ют。例：\textbf{Я зна́ю ру́сский язы́к.}
\item \textbf{ду́мать}（觉得、认为）：я ду́маю, ты ду́маешь, он ду́мает, мы ду́маем, вы ду́маете, они́ ду́мают。例：\textbf{Я ду́маю, что Даля́нь – краси́вый го́род.}
\item \textbf{… , что …}：说「知道／认为某件事」，\textbf{что 前打逗号}，从句语序不变。例：\textbf{Я зна́ю, что Минск – белору́сский го́род.}
\item \textbf{… , потому́ что …}：回答 \textbf{Почему́?}，\textbf{потому́ что 前打逗号}，跟在主句后面。例：\textbf{Мне нра́вится Даля́нь, потому́ что э́то краси́вый го́род.}
\end{itemize}

\textbf{补：这两个连接词的分工}。что 后面跟的是「内容」（知道什么、认为什么）；потому́ что 后面跟的是「原因」（为什么）。同一个句子里两个都能出现：\textbf{Я зна́ю, что … , потому́ что …}——先内容，后原因。

\subhdr{5.5　ИК-5：како́й 的第二种用法——感叹}

\S1.6 讲过 ИК-5，这里补它的用法：\textbf{како́й} 不只能提问，还能说「多么……」（У8 p27--p28）：

\begin{minipage}{\linewidth}
\tabtop
\begin{tabular}{@{}>{\raggedright\arraybackslash}p{0.16\linewidth}>{\raggedright\arraybackslash}p{0.34\linewidth}>{\raggedright\arraybackslash}p{0.34\linewidth}@{}}
\thead{句式} & \thead{句子} & \thead{中文} \\
\end{tabular}
\tabmid
\begin{tabular}{@{}>{\raggedright\arraybackslash}p{0.16\linewidth}>{\raggedright\arraybackslash}p{0.34\linewidth}>{\raggedright\arraybackslash}p{0.34\linewidth}@{}}
ИК-1 陈述 & Э́то краси́вый го́род. & 这是一座美丽的城市。 \\
ИК-3 一般疑问 & Э́то краси́вый го́род? & 这是一座美丽的城市吗？ \\
како́й 提问 & Како́й э́то го́род? & 这是座什么样的城市？ \\
ИК-5 感叹 & Како́й краси́вый го́род! & 多美的城市啊！ \\
\end{tabular}
\tabbot
\end{minipage}

\textbf{这张表在说什么}：四个句子词几乎一样，差别只在\textbf{语序（како́й 的位置）与标点}。感叹句里 како́й 放句首、句末用感叹号，调型是 ИК-5（音域拉开）。

\textbf{逐行读法}：第三行是问「什么样的」（要形容词作答）；第四行不是问，是感叹——所以不加问号。两行都出现 како́й，但一个要回答、一个不要。

\subhdr{5.6　练习：四个小题}

\begin{enumerate}[label=\textbf{第 \arabic* 题}]
\item （У8 p9）填 э́то／э́тот／э́та／э́ти：\blank мой родно́й го́род；\blank страна́ небольша́я, но краси́вая；\blank мои́ кни́ги。
\item （У8 p16）填 нра́вится／нра́вятся：（Ты）\blank белору́сская му́зыка；（Он）\blank кита́йские пельмени；（Мы）\blank кита́йский теа́тр。
\item （У8 p19）组句：я, э́ти, нра́виться, ро́зы, кра́сные。
\item （У8 p26）写问句：— Нет, не нра́вится. Мне нра́вится баскетбо́л.
\end{enumerate}

\textbf{第 1 题}　关键一步：先判名词的性数——го́род 阳、страна́ 阴、кни́ги 复。答案：э́тот ｜ э́та ｜ э́ти。

\textbf{第 2 题}　关键一步：看被喜欢的东西是单数还是复数，人一律用第三格（тебе́／ему́／нам）。答案：нра́вится ｜ нра́вятся ｜ нра́вится。

\textbf{第 3 题}　关键一步：主语是 ро́зы（复数）→ нра́вятся；кра́сные 跟 ро́зы 一致；人用 мне。答案：Мне нра́вятся э́ти кра́сные ро́зы.

\textbf{第 4 题}　关键一步：答句说「不喜欢篮球、喜欢 баскетбо́л」，问句要问「你喜欢……吗」。答案：Тебе́ нра́вится баскетбо́л?（对熟人）／Вам нра́вится баскетбо́л?（对长辈）。

\subhdr{5.7　读课文：ПАРК СИНХА́Й（У8 p29--p30）}

课件最后两页是一段完整课文，把 У8 的四个知识点全用上了（指示代词、最高级、нра́виться、ИК-5）。原文（节选，重音照课件）：

\callout{Вот парк Синха́й. Мы зна́ем, како́й э́то парк. Э́то о́чень большо́й ста́рый парк. Сейча́с о́сень. Хоро́шая тёплая пого́да. О́сенью парк о́чень краси́вый. Вот молодо́й челове́к и симпати́чная де́вушка. Они́ гуля́ют. — Мне о́чень нра́вится э́тот парк. — Мне то́же. — Посмотри́! Каки́е краси́вые ли́стья! — О́чень краси́вые: жёлтые, зелёные, кра́сные. Како́е си́нее не́бо сего́дня! — Не́бо си́нее-си́нее, как твои́ глаза́. — Како́й краси́вый парк! — Кака́я ты краси́вая!}

\textbf{课件的六问与答案}（У8 p29）：

\begin{enumerate}[label=\textbf{(\arabic*)}]
\item Како́й э́то парк? —— Э́то о́чень большо́й ста́рый парк.（两个形容词都是阳性）
\item \sloppy{}Како́й молодо́й челове́к и кака́я де́вушка гуля́ют здесь? —— Молодо́й челове́к и симпати́чная де́вушка。\par （课件只给 де́вушка 配了形容词，челове́к 没配）
\item Каки́е ли́стья о́сенью? —— Краси́вые: жёлтые, зелёные, кра́сные.
\item Како́е не́бо сего́дня? —— Си́нее.
\item Кака́я пого́да сего́дня? —— Хоро́шая тёплая.
\item Кака́я пого́да о́сенью? —— 课文里「今天」就是秋天，所以仍答 хоро́шая тёплая пого́да；课件没说「秋天总是好天气」，不能扩大。
\end{enumerate}

\textbf{关键一步}：六个问句的疑问词（како́й／кака́я／како́е／каки́е）已经给出答案的性数——答题时只做两件事：\textbf{从课文里找对应的名词}，\textbf{把形容词词尾抄成疑问词的那个族}。

\textbf{四个感叹句各用了什么}（与 5.5 对照）：\textbf{Каки́е краси́вые ли́стья!}（复数，ИК-5）、\textbf{Како́е си́нее не́бо!}（中性）、\textbf{Како́й краси́вый парк!}（阳性）、\textbf{Кака́я ты краси́вая!}（阴性）——四句四个性数，正是 5.1 那张表在真实句子里的样子。
